\chapter{Sjølvbedraget}

\begin{motto}
Ein jury som ikkje kan grunngje dommen sin, deler likevel ut prisar.\\ Det er sjølve poenget med prisen.
\end{motto}

Eit fag utan fundament burde rakne. At designfaget ikkje har rakna, at det tvert om veks, fyller universitet, deler ut grader, vinn prestisje, krev ei forklaring. Forklaringa er at faget har bygd eit apparat som har til funksjon å produsere teikn på at standardar finst, utan at standardane treng å finnast. Akkrediteringa, prisane, fagfellevurderinga: kvar av dei ser ut som ein mekanisme for kvalitetssikring, og kvar av dei er i røynda ein mekanisme for å sirkulere prestisje innanfor ein lukka krins. Til saman utgjer dei sjølvbedraget, ikkje løgn frå nokon einskild, men ein kollektiv struktur som gjer at ingen treng å sjå holet.

\section{Akkrediteringa måler det målbare}

Sjå fyrst på korleis ei designutdanning vert godkjend. Eit panel kjem på besøk, granskar pensum, ser på fasilitetar, intervjuar tilsette og studentar, og godkjenner. Men sjå kva panelet faktisk kan vurdere. Det kan telje at det finst nok lærarar med rett formell kompetanse. Det kan stadfeste at det finst verkstader og programvare. Det kan kontrollere at læringsutbyteformuleringane er på plass. Med andre ord: akkrediteringa måler det målbare, ressursar, prosessar, dokument, og må ta det umålbare, sjølve kvaliteten på det studentane lærer å sjå og gjere, på tru. Ho kan ikkje gjere anna, for å måle det umålbare ville krevje den standarden faget ikkje har.

Det lønner seg å gå tett på eit konkret organ, for abstraksjonen «akkreditering» løyner kor presist mekanismen omgår spørsmålet om kvalitet. Ta National Association of Schools of Art and Design, NASAD, skipa i USA i 1944 og framleis den instansen som godkjenner dei fleste kunst- og designutdanningane der. NASAD fører ei handbok med det dei kallar «standards», og det lønner seg å lese kva desse standardane faktisk gjeld \autocite{nasad_handbook}. Dei gjeld talet på krediteringstimar i eit program. Dei gjeld dei formelle kvalifikasjonane til lærarstaben. Dei gjeld biblioteket, verkstaden, utstyret, arealet per student. Dei gjeld «health and safety» i verkstadene. Dei gjeld at studieplanen finst som dokument, at læringsutbyta er nedskrivne, at det finst rutinar for vurdering. Med eitt ord: standardane gjeld \emph{ressursar} og \emph{prosess}.

Sjå kva som ikkje står der, og ikkje kan stå der. Ingen stad i handboka finst ein målestokk for sjølve den estetiske eller formmessige dømekrafta studentane skal lære, for ein slik målestokk ville krevje den standarden faget ikkje har. NASAD kan stadfeste at ein skule \emph{har} det som skal til for å drive same verksemd som alle andre godkjende program. Han kan ikkje seie at det som blir lært, er godt, fordi «godt» her ikkje har noko utanfor krinsen å bli målt mot. Akkrediteringa flytter altså spørsmålet eitt hakk til sida: frå «er det studentane lærer å sjå og gjere, godt?» til «har skulen ressursane til å lære dei det same uvisse som dei andre lærer?» Det fyrste spørsmålet kan ho ikkje svare på; det andre svarar ho grundig på. Og fordi det andre ser ut som det fyrste, fordi ein godkjend skule kjennest som ein god skule, treng ingen merke byttet.

Prosedyren seier kva slag instans dette er. Akkreditering hjå NASAD går føre seg ved at skulen fyrst lagar ein omfattande «self-study», ei sjølvgransking mot standardane, og deretter får besøk av eit «visiting evaluation team» sett saman av jamlikar frå \emph{andre} akkrediterte program. Det er den same lukka krinsen som vurderer seg sjølv. Dei som allereie er innanfor, avgjer kven som slepp inn, etter kriterium dei sjølve forvaltar, granska av folk frå hus som er godkjende etter dei same kriteria. Sirkelen er ikkje skjult; han er sjølve metoden. Difor er ho streng der ho kan telje og taus der det gjeld: ein skule fell på manglande areal, uklåre studieplandokument eller for få kvalifiserte lærarar, aldri på at det studentane lærer å sjå, er feil, for det finst ingen standard å felle den dommen frå. Resultatet av at alle godkjende program er målte mot dei same ressurs- og prosesskrava, er ikkje at programma blir gode. Det er at dei blir like.

At Alain Findeli har etterlyst eit klarare grunnlag for nett denne utdanninga \autocite{findeli2001}, og at etterlysinga ikkje har endra akkrediteringspraksisen, fortel kor djupt skinet er forankra. Ei innanfrå-kritikk prellar av, ikkje fordi nokon avviser ho, men fordi ho ikkje har noko å gripe i. Apparatet held fram med å gjere det same det alltid har gjort, og kallar det kvalitetssikring.

\section{Den gule, den kvite, den svarte blyanten}

Prisane er endå tydelegare, og dei to fremste i faget syner kvar sin variant av den same sirkelen. Sjå fyrst på D\&AD, British Design \& Art Direction, skipa i London i 1962 av ein flokk fotografar og designarar, mellom dei David Bailey, Terence Donovan, Alan Fletcher og Colin Forbes, med dei fyrste prisane utdelte året etter. D\&AD deler ut det dei kallar blyantar, Yellow Pencil, White Pencil, og på toppen Black Pencil, og det ber alle teikn på ein streng standard. Black Pencil blir somme år rett og slett \emph{ikkje} delt ut i ein kategori: juryen held att, fordi ingenting var godt nok. Her, kunne ein tru, har vi endeleg ein målestokk med tenner, ein som tør å seie nei.

Men sjå kva tilbakehaldet faktisk er. Det er framleis juryens delte smak som avgjer at ingenting heldt mål, ikkje ein ekstern standard som arbeida fell på. At ein jury samstemmer om at ingen fortener prisen i år, syner at juryen har ein høg terskel; det syner ikkje at terskelen er noko anna enn deira eige, ikkje-overførbare blikk. Tilbakehaldet er det næraste faget kjem eit nei, og nett difor er det avslørande: når faget endeleg seier nei, kviler nei-et på det same blikket som elles seier ja, og ikkje på noko utanfor blikket. Og juryane er sette saman av bransjefolk, dei etablerte, som kårar dei lovande, som dermed blir dei nye etablerte og sit i neste jury. Tilbakehaldet av ein pris er ikkje strengskapens prov; det er smakens sjølvtillit, kledd som måling.

Compasso d'Oro syner den andre varianten: ikkje strengskap, men kanondaning. Prisen vart skipa i Italia i 1954, på ein idé frå Gio Ponti og på initiativ frå varehuset La Rinascente, som dreiv dei fyrste fire utgåvene sjølv før Associazione per il Disegno Industriale, ADI, tok over. Det er den eldste industridesignprisen som finst, og vinnararbeida går inn i ei «Historical Collection» som no tel over to tusen tre hundre objekt og har sitt eige hus, ADI Design Museum i Milano, opna i 2021. Her ser ein mekanismen reint: prisen \emph{konstituerer} kanon. Det framifrå er det krinsen kårar, og det krinsen kårar, blir samla, stilt ut og overlevert til ettertida som det framifrå. Ein seinare granskar som spør «kva var god italiensk design?» blir vist til samlinga, som er nett dei arbeida krinsen alt hadde kåra.

Lukkinga blir kalla arv, og det er ikkje berre ein talemåte. I 2004 erklærte det italienske kulturdepartementet samlinga for å vere av eksepsjonell kunstnarleg og historisk interesse, eit nasjonalt kulturminne verna ved lov. Staten har då formelt stadfest at smaken til krinsen er kulturarv. Dommen krinsen felte om seg sjølv, har fått eit segl frå utsida som ser ut som ei uavhengig stadfesting og er ei innramming av den same dommen.

Bak begge prisane ligg det same spørsmålet boka har stilt heile vegen: kviler prisen på eit blikk som kan overførast, eller på eit blikk som berre kan delast av dei som alt har det? Prøv prøven, ein gong, skikkeleg. Kunne ein uavhengig jury, utan kontakt med den fyrste, kome til same dom \emph{av same grunn}? For prisen, som for kritikken i atelieret, er svaret nei. Ein annan jury kunne tilfeldigvis kåre same arbeid, men ikkje av ein grunn den fyrste juryen kunne ha skrive ned og overlevert; dei kunne berre kome dit ved sjølve å ha det same trente blikket, det vil seie ved alt å høyre til krinsen. Prisen måler ikkje kvalitet mot ein standard. Han sirkulerer prestisje innanfor ein smak, og festar smaken til veggen i eit museum så han ser ut som kunnskap.

\section{To drakter på same gest}

Prøven gjev same svar kvar gong, og det syner seg klarast når ein set dei to mest ulike scenene i faget ved sida av kvarandre. Den eine er offentleg og kledd i ljos; den andre er privat og kledd i frist. Under begge ligg den same gesten.

Tenk deg fyrst prisutdelinga av det store slaget. Ein scene i mørklagd sal, runde bord med duk og ljos, ein konferansier som les opp nominasjonane medan ein storskjerm syner arbeida i tur. Tre hundre menneske i mørket ventar på namnet i konvolutten. Juryen, sju, åtte fagfolk, sit allereie ute i salen, ferdige med arbeidet sitt; dommen er felt før gjestene kom. Tidlegare same dagen sat dei same åtte med bidraga framfor seg. Dei hadde ikkje rekna; dei hadde \emph{sett}. Ein hadde sagt at dette arbeidet «har noko», at det «sit», at det er «modig der dei andre er trygge». Dei andre nikka, for dei såg det same, og det er nett samstemma som kjennest som prov. Bytt ut dei åtte med åtte andre frå same miljø, og dei kunne kåre eit anna arbeid og grunngjeve det like varmt. Be dei skrive ned grunngjevinga så presist at ein framand kunne velje rett utan å sjå arbeida, og dei ville ikkje klare det. Grunngjevinga finst, ho er velformulert, ho nemner mot og klårleik og relevans, men ho er ei \emph{omskriving} av dommen, ikkje ein veg til han.

Flytt no scena. Ein person åleine framfor ein skjerm, klokka over midnatt, med eit manuskript dei har sagt ja til å fagfellevurdere. Ingen sal, ingen ljos, ingen klapp, berre eit skjema med felt som skal fyllast ut og ein frist som nærmar seg. Manuskriptet er på tjuefem sider, velskrive, plasserer seg pent i litteraturen, siterer dei rette namna, har ein metodedel som lyder rett. Fagfellen les, og medan han les, formar det seg ein dom, ikkje av eit kriterium, men av ei \emph{kjensle}: dette held mål. Dette er av det slaget vi publiserer. Spør kva den kjensla målte. Ikkje om påstanden er sann, om resultatet held, om nokon som gjorde det same ville finne det same. Kjensla målte noko meir beskjedent og meir avslørande: om arbeidet \emph{liknar} det krinsen plar godkjenne. «Det kjennest seriøst» er ein dom om sjanger, ikkje om sanning, og i eit fag utan ein delt målestokk er sjanger alt fagfellen har.

Det er ein og same gest, sett frå to kantar. Eit trent blikk frå ein krins møter eit arbeid, kjenner at det «sit» eller «held mål», og feller ein dom som ser ut som ei måling og er ei attkjenning. Kunne ein uavhengig fagfelle kome til same dom av same grunn? Nei, av nett same grunn som juryen ikkje kunne. Det skjuler holet er drakta: i salen ritualet, ljoset, klappinga, trofeet i ljoset; framfor skjermen prosedyren, skjemaet, fristen, ordet «fagfellevurdert» trykt på artikkelen etterpå. Apparatet produserer i begge tilfelle eit teikn på at ein standard har vore i verksemd, utan at standarden treng å ha funnest.

Det er freistande å tenkje at forskinga i det minste byggjer seg opp over tid, sjølv om gallaen berre feirar. Christopher Frayling forsøkte å gje designforskinga eit grunnlag ved å dele ho i forsking inn i, gjennom og for design \autocite{frayling1993}, ei inndeling som har vorte sitert utan ende. Men inndelinga ordnar slag av forsking; ho gjev ikkje eit felles kriterium for kva som tel som godt arbeid innanfor noko av slaga. Difor konvergerer ingenting. Ein artikkel motbeviser ikkje ein annan; han legg seg ved sida av. Bruce Archer kravde at design skulle vere ein eigen disiplin med eige kunnskapsgrunnlag \autocite{archer1979discipline}, og Nigel Cross forsvarte den posisjonen så godt ho lèt seg forsvare ved å peike på «designer ways of knowing», ein eigen, ikkje-verbal dømekraft som skil designaren frå vitskapsmannen \autocite{cross2006}. Forsvaret er skarpt, men det redder ikkje apparatet; det \emph{skildrar} det. Cross gjev oss eit presist namn på blikket fagfellen feller dommen med, og syner i same andedrag kvifor dommen ikkje let seg overlevere: ein dømekraft som verkar nett ved å ikkje vere verbal, kan ikkje skrivast ned og sendast til ein framand. Det fagfellevurderinga produserer over tid, er difor ikkje ein kumulativ kropp, men eit arkiv av meiningar med fotnotar, sortert etter kor godt dei kjennest. Drakta forskinga ber, metodedelen, litteraturoversynet, fagfellevurderinga, er lånt frå den kumulative kunnskapsbygginga. Kritikken er ikkje at designforskinga er ein mislukka vitskap, for ho er ikkje eit forsøk på å vere det same slaget som vitskapen. Kritikken er at ho ber drakta til noko ho ikkje er, og let drakta gjere arbeidet grunnen skulle ha gjort.

\section{Kvifor alle spelar med}

Det som held sjølvbedraget oppe er ikkje vondskap eller dumskap, men interesse. Skulane treng studentar, og studentar treng grader som verkar verdifulle. Studentane har betalt og investert år, og treng at grada tyder noko. Lærarane har bygd karrierar i faget, og treng at faget er reelt. Arbeidsgjevarane treng eit signal å tilsetje etter, og grada er signalet dei har. Kvar aktør, rasjonelt, vil halde oppe trua på at her finst eit fag med standardar, fordi alternativet, å seie høgt at her er ingen standard å måle mot, ville koste kvar av dei noko. Sjølvbedraget er ein likevekt, ikkje ein konspirasjon. Ingen lyg. Alle deltek i ein struktur som produserer teikn på kvalitet, og alle har grunn til å lese teikna som ekte. Det er nett difor det er så stabilt, og nett difor det ikkje kjem til å rive seg sjølv ned innanfrå. Likevekter held til noko utanfrå bryt dei.

\vignett

Og noko utanfrå er på veg. Eit fag kan halde oppe eit sjølvbedrag så lenge ingen ekstern kraft krev av det noko det ikkje kan levere. Den krafta kjem no frå to kantar samstundes: frå teknologiar som tvingar fram presise spørsmål faget ikkje kan svare på, og frå historia, som har sett dette mønsteret før og veit korleis det endar.

\laeringsmal{\begin{itemize}
\item Gjere greie for skilnaden mellom ein ressurs- og prosesstandard (som NASAD-handboka set) og ein målestokk for sjølve kvaliteten på det som blir lært, og forklare kvifor akkrediteringa berre kan femne den fyrste.
\item Forklare korleis «self-study» og «visiting evaluation team» gjer akkrediteringa til ein krins som vurderer seg sjølv, og kva som ville krevjast for at ho ikkje var det.
\item Skildre korleis ein designpris både kan synast streng (D\&AD som held att Black Pencil) og konstituere kanon (Compasso d'Oro og samlinga, verna som kulturminne), og kvifor ingen av delane er ei måling mot ein ekstern standard.
\item Bruke prøven «kunne ein uavhengig jury kome til same dom av same grunn?» på akkreditering, pris og fagfellevurdering, og kjenne att kvifor svaret er det same kvar gong.
\end{itemize}}

\ovingar{\begin{enumerate}
\item Finn fram standardane til eit akkrediteringsorgan for design (NASAD eller eit nasjonalt motstykke). List opp kva kvar standard faktisk måler. Marker kvar einaste der ein ressurs eller prosess står i staden for ein dom over kvalitet. Kor mange att er det av det siste slaget?
\item Vel ein vinnar av Compasso d'Oro eller ein Black Pencil. Skriv den beste grunngjevinga du kan for dommen. Gjev så teksten til ein som ikkje har sett arbeidet, og spør om grunngjevinga åleine ville late dei kåre same arbeid framfor eit anna. Kva manglar?
\item Eit år blir Black Pencil ikkje delt ut i ein kategori. Drøft: er dette prov på ein streng standard, eller på ein sjølvsikker smak? Kva slags evidens skulle til for å skilje dei to?
\item Designakkrediteringa kan ikkje måle kvaliteten på det studentane lærer å sjå. Akkrediteringa av eit medisinprogram kan derimot peike på pasientutfall utanfor krinsen. Drøft kva slags fundament dei to kviler på, og kvifor fråvêret av eit slikt fundament i designfaget ikkje gjer det til ei mislukka utgåve av medisinen, men til noko av eit anna slag som må grunngjevast på sine eigne premiss.
\end{enumerate}}

\vidarelesing{\fullcite{archer1979discipline}; \fullcite{cross2006}; \fullcite{findeli2001}; \fullcite{frayling1993}.}
